\chapter{HTTP that matters}

HTTP is the protocol your backend speaks.
You do not need every feature, but you need the fundamentals.

\section{Methods and semantics}
Use methods for intent:
\begin{itemize}[nosep]
  \item \code{GET} for reads
  \item \code{POST} for creation
  \item \code{PATCH} for partial updates
  \item \code{DELETE} for removal
\end{itemize}

\section{Status codes}
Status codes are signals to clients and monitoring.
Use them consistently.

\begin{itemize}[nosep]
  \item \code{200} success
  \item \code{201} created
  \item \code{400} validation error
  \item \code{401} unauthenticated
  \item \code{403} unauthorized
  \item \code{404} not found
  \item \code{409} conflict
  \item \code{500} unexpected error
\end{itemize}

\section{Status code mapping}
Write down a mapping table.
Keep it consistent across endpoints:

\begin{tabularx}{\textwidth}{@{}lX@{}}
\toprule
Code & Typical meaning \\
\midrule
200 & Successful read or update \\
201 & Resource created successfully \\
202 & Accepted for background processing \\
204 & Successful response with no body \\
400 & Validation or parsing error \\
401 & Authentication missing or invalid \\
403 & Authenticated but not allowed \\
404 & Resource not found \\
409 & Conflict with existing data \\
429 & Rate limit exceeded \\
500 & Unhandled server error \\
\bottomrule
\end{tabularx}

\section{Idempotency}
Writes should be safe to retry.
Use idempotency keys for payment and job creation endpoints.

\section{Headers that matter}
Headers carry important semantics:
\begin{itemize}[nosep]
  \item \code{Content-Type} for payload format
  \item \code{Accept} for response negotiation
  \item \code{Authorization} for authentication
  \item \code{Idempotency-Key} for safe retries
\end{itemize}

\section{Idempotency in practice}
Store the request hash and response for a short period.
If the same key appears again, return the stored response.

\begin{lstlisting}[language=JSON]
{
  "idempotency_key": "req_9c6",
  "status": "stored",
  "expires_in_seconds": 86400
}
\end{lstlisting}

\section{Pagination}
Never return unbounded lists.
Choose cursor-based pagination for large datasets.

\begin{lstlisting}[language=JSON]
{
  "items": [
    {"id": "task_123", "title": "Write docs"}
  ],
  "next_cursor": "task_123"
}
\end{lstlisting}

\section{Caching headers}
Set caching headers for read-heavy endpoints.
Use \code{ETag} or \code{Last-Modified} to reduce bandwidth.

\section{Content negotiation}
Support at least JSON and make it explicit.
Reject unexpected content types with a clear error.

\section{Safe and unsafe methods}
Clients and proxies treat methods differently.
Ensure \code{GET} is safe and idempotent, and use \code{POST} for non-idempotent writes.

\section{Streaming responses}
If results are large, consider streaming or chunked responses.
This keeps memory usage stable and improves time to first byte.

\section{Request size limits}
Set request size limits to protect memory and avoid abuse.
Document those limits for clients.

\section{Problem details}
Consider using a standard error envelope or problem details format.
Consistency is more important than the specific shape.

\section{Error response example}
\begin{lstlisting}[language=JSON]
{
  "error": {
    "code": "validation_error",
    "message": "title is required",
    "field": "title"
  }
}
\end{lstlisting}

\section{Request body size}
Set a maximum request body size.
Large bodies can exhaust memory and slow downstream systems.

\section{Client retry behavior}
Document which endpoints are safe to retry.
If a client retries non-idempotent requests, you will get duplicate data.

\section{Caching strategy}
Use \code{Cache-Control} to define shared and private caching.
If responses contain user-specific data, avoid shared caches.

\section{Content type versioning}
If you need versioning, consider vendor media types.
It is more explicit than path-based versioning and does not change URLs.

\section{Conditional requests}
Use conditional requests to avoid sending full responses when data has not changed.
If the client presents an \code{If-None-Match} value and the resource is unchanged,
return \code{304 Not Modified}.

\section{ETag example}
\begin{lstlisting}[language=JSON]
{
  "etag": "task_123:v7",
  "last_modified": "2026-02-18T10:00:00Z"
}
\end{lstlisting}

\section{Rate limiting}
Expose rate limits via headers so clients can behave responsibly:
\begin{itemize}[nosep]
  \item \code{X-RateLimit-Limit}
  \item \code{X-RateLimit-Remaining}
  \item \code{X-RateLimit-Reset}
\end{itemize}

\section{Retry-After}
When returning \code{429} or \code{503}, include \code{Retry-After}.
This helps clients back off and prevents thundering herds.

\section{Timeouts}
Clients should not wait forever.
Publish expected timeouts for critical endpoints and enforce them on the server.

\section{Compression}
Enable gzip or brotli for large responses.
Compression reduces bandwidth and often improves tail latency.

\section{Correlation identifiers}
Include a request id in every response and log.
This enables fast incident debugging and request tracing.

\section{Example request}
\begin{lstlisting}[language=Bash]
curl -i https://api.example.com/tasks \
  -H "Authorization: Bearer <token>" \
  -H "Accept: application/json"
\end{lstlisting}

\section{Example response}
\begin{lstlisting}[language=JSON]
{
  "items": [
    {"id": "task_123", "title": "Write docs"}
  ],
  "next_cursor": "task_123"
}
\end{lstlisting}

\section{Health and readiness endpoints}
Expose simple health checks for load balancers and deployment systems.
Keep them fast and side-effect free.

\section{Batch endpoints}
Batch endpoints can reduce request overhead but increase complexity.
Use them only when the latency or cost savings are meaningful.

\section{CORS}
If browsers call your API directly, configure CORS explicitly.
Avoid broad wildcards for production systems.

\section{Security headers}
Security headers like \code{Strict-Transport-Security} and \code{X-Content-Type-Options}
reduce attack surface for browser-based clients.

\section{Authentication patterns}
Be explicit about how clients authenticate:
\begin{itemize}[nosep]
  \item short-lived bearer tokens for user-facing clients
  \item long-lived API keys for server-to-server calls
  \item mTLS for highly sensitive internal services
\end{itemize}

Each pattern implies different rotation and monitoring requirements.

\section{Authorization errors}
Distinguish between unauthenticated and unauthorized.
Return \code{401} when authentication is missing or invalid, and \code{403} when the user is authenticated but not allowed.
This improves client behavior and audit clarity.

\section{Request validation}
Validate at the edge and return structured errors.
Reject unknown fields if you want a strict contract.
If you allow unknown fields, document that explicitly and ignore them safely.

\section{Partial responses}
Large responses can be expensive to produce and transmit.
Allow clients to request a subset of fields, and default to a safe, stable set.
This keeps the API responsive as resources grow.

\section{Webhooks and callbacks}
Webhooks are part of the HTTP surface area.
Treat them as public endpoints:
\begin{itemize}[nosep]
  \item sign webhook payloads
  \item include a request id for retries
  \item document the retry strategy and time window
\end{itemize}

\section{Long-running work}
If work takes more than a few seconds, return \code{202 Accepted} with a job id.
Expose a \code{GET /jobs/:id} endpoint so clients can poll.

\section{Error budget for clients}
Client behavior is part of reliability.
Publish clear guidance:
\begin{itemize}[nosep]
  \item which endpoints are safe to retry
  \item the maximum recommended retry window
  \item how to detect throttling
\end{itemize}

\section{Input size and pagination limits}
Define maximum page size and request body size.
When a client exceeds limits, respond with \code{400} and explain the constraint.

\section{Conditional request example}
Clients can avoid full responses using conditional requests:
\begin{lstlisting}[language=Bash]
curl -i https://api.example.com/tasks/task_123 \
  -H "If-None-Match: \"task_123:v7\""
\end{lstlisting}

If the resource is unchanged, the server returns \code{304} with no body.

\section{ETag and If-Match for writes}
Use \code{If-Match} to prevent lost updates.
If a client sends a stale ETag, return \code{412 Precondition Failed}.
This is a clean way to protect concurrent updates.

\section{File uploads}
If you accept files, document limits and formats.
Large uploads should use multipart uploads or pre-signed URLs to keep the API responsive.

\section{HTTP caching workflow}
For read-heavy endpoints, caching should be part of the workflow:
\begin{itemize}[nosep]
  \item server returns \code{ETag} and \code{Cache-Control}
  \item client revalidates with \code{If-None-Match}
  \item server returns \code{304} if unchanged
\end{itemize}

This reduces bandwidth and avoids unnecessary processing.

\section{Problem details with metadata}
If you use a problem details format, include a stable error code and a trace id.
This makes it easier to correlate client failures with server logs.

\section{Header budget}
Headers are not free.
Large headers increase latency and can exceed gateway limits.
Keep custom headers minimal and document those that matter.

\section{Response shape stability}
Clients depend on response shape.
Avoid changing data types or omitting fields for existing endpoints.
If you need a new representation, introduce a new field or a new endpoint.

\section{Cache-Control directives}
Make caching explicit with \code{Cache-Control}:
\begin{itemize}[nosep]
  \item \code{max-age} for browser and client caches
  \item \code{s-maxage} for shared caches
  \item \code{private} for user-specific responses
\end{itemize}

\section{HTTP response example}
An explicit response makes behavior visible to clients:
\begin{lstlisting}[language=Bash]
HTTP/1.1 200 OK
Content-Type: application/json
Cache-Control: private, max-age=30
ETag: "task_123:v7"
X-Request-Id: req_9c6

{"id":"task_123","title":"Write docs","status":"open"}
\end{lstlisting}

\section{Cookies and sessions}
If you use cookies, set \code{HttpOnly} and \code{Secure} flags.
Consider \code{SameSite=Lax} or \code{Strict} to reduce CSRF exposure.
Even token-based APIs should document how cookies are handled.

\section{Redirects}
Avoid redirects for API endpoints.
They confuse clients and complicate retries.
If you must redirect, use \code{307} or \code{308} to preserve method semantics.

\section{Method safety matrix}
Document which methods are safe and idempotent.
This guides client retries and caching:
\begin{tabularx}{\textwidth}{@{}lXX@{}}
\toprule
Method & Safe & Idempotent \\
\midrule
GET & yes & yes \\
POST & no & no \\
PUT & no & yes \\
PATCH & no & depends on semantics \\
DELETE & no & yes (if defined that way) \\
\bottomrule
\end{tabularx}

\section{Content negotiation example}
If you support multiple formats, show a concrete example:
\begin{lstlisting}[language=Bash]
curl -i https://api.example.com/tasks/task_123 \
  -H "Accept: application/json"
\end{lstlisting}

If a client requests an unsupported format, return \code{406 Not Acceptable}.

\section{Link headers for pagination}
Link headers can make pagination easier for clients:
\begin{lstlisting}[language=Bash]
Link: <https://api.example.com/tasks?cursor=task_123>; rel="next"
\end{lstlisting}

If you use Link headers, keep them consistent with the response body.

\section{Range requests}
For large downloads, support \code{Range} when it makes sense.
This allows clients to resume transfers without re-downloading the entire file.

\section{HTTP error response example}
When errors happen, include headers and a structured body:
\begin{lstlisting}[language=Bash]
HTTP/1.1 422 Unprocessable Entity
Content-Type: application/json
X-Request-Id: req_42

{"error":{"code":"validation_error","message":"title is required","field":"title"}}
\end{lstlisting}

\section{Throttling strategy}
Rate limits should be predictable.
Document the limits per endpoint or per client type.
If you use dynamic limits, return them via headers so clients can adapt.

\section{Batch request example}
Batch endpoints should document partial failures clearly.
An example batch response:
\begin{lstlisting}[language=JSON]
{
  "results": [
    {"id": "task_1", "status": "created"},
    {"id": null, "status": "error", "error": "validation_error"}
  ]
}
\end{lstlisting}

If you allow batch requests, define how you handle retries for partial success.

\begin{checklistbox}{Checklist: HTTP}
\begin{itemize}[nosep]
  \item Use methods for intent.
  \item Pick status codes and stick to them.
  \item Add pagination before you need it.
\end{itemize}
\end{checklistbox}
